\subsection{正交性}

定义3. 设 E 是一个 A-模， \(\mathbf{E}^*\) 为其对偶；一个元素 \(x \in \mathbf{E}\) 与一个元素 \(x^* \in \mathbf{E}^*\) 称为正交的，如果 \(\langle x, x^* \rangle = 0\) .

E 的子集 M 与子集 \(M'\) （属于 \(E^{*}\) ）称为正交的，如果对所有 \(x \in M\) , \(x^{*} \in M'\) ，x 与 \(x^{*}\) 正交。特别地， \(x^{*} \in E^{*}\) （相应地 \(x \in E\) ) 称为与 M （相应地 \(M'\) ) 正交，如果它与 M （相应地 \(M'\) ) 的每个元素正交。如果 \(x^{*}\) 和 \(y^{*}\) 与 M 正交，那么 \(x^{*} + y^{*}\) 和 \(x^{*}\alpha\) 对所有 \(\alpha \in A\) 也与 M 正交（由第3号的(10)和(12)），这就证明了以下定义的合理性：

定义4. 给定E的一个子集M（相应地，E*的一个子集M'），与M（相应地，M'）正交的 \(x^* \in \mathbf{E}^*\) 的集合（相应地， \(x \in \mathbf{E}\) 的集合）称为与M（相应地，M'）完全正交的子模（或者，若不会引起混淆，简称为与M（相应地，M'）正交的子模）。

根据线性形式的定义， \(E^{*}\) 中与E正交的子模简化为0； \(E^{*}\) 中与 \(\{0\}\) 正交的子模恒等于 \(E^{*}\) .

命题6. 设 M, N 是 E 的两个子集，且 M ⊂ N；若 M' 与 N' 分别是 E* 中正交于 M 与 N 的子模，则 N' ⊂ M'。

命题7. 设 \((\mathbf{M}_t)\) 是E的一族子集；与诸 \(\mathbf{M}_t\) 之并正交的子模是诸子模 \(\mathbf{M}_t'\) 的交，这些子模分别正交于对应的 \(\mathbf{M}_t\) ；这个子模也是与E中由诸 \(\mathbf{M}_t\) 之并生成的子模正交的子模。

这些结果是定义的直接推论。

对于E中与 \(E^{*}\) 的子集正交的子模，存在一个类似的命题（其表述留给读者）。

若M是E的一个子模， \(M'\) 是 \(E^{*}\) 中与 M 正交的子模，且 \(M''\) 是E中与 \(M'\) 正交的子模，则 \(M \subset M''\) 但可能有 \(M \neq M''\) （习题9）。然而请注意，若 \(M''\) 是 \(M''\) 在 \(E^{*}\) 中的正交，则 \(M'' = M'\) ；因为 \(M' \subset M''\) ，而另一方面，关系 \(M \subset M''\) 蕴含 \(M'' \subset M'\) .

